\documentclass[11pt]{article}
\usepackage[margin=1in]{geometry}
\usepackage[T1]{fontenc}
\usepackage[utf8]{inputenc}
\usepackage{lmodern}
\usepackage{hyperref}
\usepackage{xcolor}
\usepackage{fancyvrb}
\usepackage{listings}
\usepackage{CJKutf8}
\lstset{basicstyle=\ttfamily\small,columns=fullflexible,breaklines=true,upquote=true,frame=single,framerule=0.2pt,tabsize=2}

\title{\texttt{\detokenize{model/__init__.py}} (Q\&A Documentation)}
\author{}
\date{}
\begin{document}
\begin{CJK*}{UTF8}{gbsn}
\maketitle

\paragraph{Q: What is the purpose of this package initializer?}
\textbf{A:} It exposes a convenient import surface for the package, so scripts can import key classes from the package root.

\paragraph{Q: Why is this helpful for the project?}
\textbf{A:} It stabilizes imports and keeps training scripts cleaner and less coupled to internal file layout.

\paragraph{Q: What does this file export (public API)?}
\textbf{A:} The public API is the set of classes/functions other modules import (sometimes re-exported by package initializers). The key top-level definitions detected here are:

\begin{Verbatim}[fontsize=\small]
# (no top-level classes/functions detected)
\end{Verbatim}

\paragraph{Q: What are the main dependencies (imports) and why do they matter?}
\textbf{A:} Imports show whether this file is model code (torch), data code (pandas), or evaluation/analysis (sklearn). They also determine runtime requirements.

\vspace{1mm}
\VerbatimInput[firstline=1,lastline=6,fontsize=\small]{/workspace/model/__init__.py}


\paragraph{Q: Example: end-to-end data flow involving this file}
\textbf{A:} Core pipeline shape (schematic):

\begin{Verbatim}[fontsize=\small]
for features, label in dataloaders["train"]:
    out = model(features)
    loss = ...
    loss.backward(); optimizer.step()
\end{Verbatim}

\paragraph{Q: Full source code of the Python file}
\textbf{A:} The complete file is included verbatim for reference.

\VerbatimInput[fontsize=\scriptsize]{/workspace/model/__init__.py}

\end{CJK*}
\end{document}